\section{Preliminaries}
\label{sec:preliminaries}

\paragraph{Reinforcement Learning for LLMs}

Reinforcement Learning (RL)~\citep{Sutton1988ReinforcementLA} is general paradigm for sequential decision-making problems. RL has demonstrated its great ability of optimizing pre-trained LLMs~\citep{trung2024reft,jaech2024openai}.
In this paradigm, sequential token generation is modeled as a Markov Decision Process (MDP) $M = (\mathcal{S}, \mathcal{A}, P, R, \gamma)$,
where a \textbf{state} $s_t = (q, y_{1:t}) \in \mathcal{S}$ is the prompt with the output generated so far, and the \textbf{action} $a_t \in \mathcal{A}$ is the next token selected from the vocabulary $\mathcal{V}$. 
Hence, the \textbf{transition $P(s_{t+1}|s_t, a_t)$} is deterministic in this context.
An episode starts from a prompt $s_0$ (out of a predefined set $\mathcal{D}_0$) and terminates at an end-of-sequence token or by the maximum sequence length $H$.
The \textbf{reward $R(s_t, a_t)$} signal is issued by either a rule-based reward function or a learned reward model.
In the scope of this paper, we consider the verifiable reward function. For any non-terminal timestep $t < T-1$, $R(s_t, a_t)$ is $0$; on completion, the terminal reward, denoted by $R(\tau)$ for the whole sequence, equals $1$ if $\tau$ produces a correct and well-formatted answer and $0$ otherwise.

The \textbf{policy $\pi_\theta(a_t \mid s_t)$} in the MDP is the LLM itself, parameterized by $\theta$, and it defines a probability distribution of next-token generation. 
We use $d^{\pi_{\theta}}_{\tau}$ to denote the distribution of the output sequence $\tau$ generated by $\pi_{\theta}$ and use $d^{\pi_{\theta}}_{s,a}, d^{\pi_{\theta}}_{s}$ for the state-action pairs $(s,a)$ and the state respectively. 
The learning objective of an RL policy is to maximize the reward function, i.e., $\pi^* = \arg \max_{\pi_{\theta}} J(\pi_{\theta})$, where $J(\pi)
    =
    \mathbb{E}_{\tau \sim \pi}
    \left[
        R(\tau)
    \right]$. 


\vspace{-0.2cm}

\paragraph{Monotonic Policy Improvement and Proximal Policy Gradient}

Optimizing RL policy according to the vanilla policy gradient $\nabla_{\pi}J(\pi)$ often turns out to be ineffective as the performance is sensitive to the step size. To this end, \citet{SchulmanLAJM15TRPO} revisits the \textit{policy difference}~\citep{KakadeL02CPI} for arbitrary two policies $\pi$ and $\pi^{\prime}$ as defined below:
\begin{equation}
    \Delta(\pi^{\prime},\pi)
    :=
    J(\pi^{\prime}) - J(\pi)
    =
    \mathbb{E}_{s,a \sim d^{\pi^{\prime}}}
    \left[
        A^{\pi}(s,a)
    \right],
    \label{eq:exact_policy_difference}
\end{equation}
where $d^{\pi^{\prime}}$ denotes the discounted state visitation distribution induced by policy $\pi^{\prime}$, and let $A^{\pi}(s,a)$ denote the advantage of taking action $a$ at state $s$ under policy $\pi$. 
Then, a natural idea is to ensure \textit{monotonic policy improvement} during iterative policy optimization.
Formally, given an old policy $\pi_k$, the goal of monotonic policy improvement is to find a new policy $\pi_{k+1}$ such that $J(\pi_{k+1})-J(\pi_k)\geq 0$, i.e., $\Delta(\pi_{k+1}, \pi_k) \ge 0$.



However, directly optimizing~\autoref{eq:exact_policy_difference} is infeasible because it depends on the visitation distribution of the updated policy, i.e., $\pi_{k+1}$. 
\citet{SchulmanLAJM15TRPO} propose Trust-Region Policy Optimization (TRPO) to approximate the exact objective difference by replacing the updated visitation distribution with the old-policy distribution and estimating with the importance sampling ratio:
\begin{equation}
    \widetilde{\Delta}(\pi_{k+1},\pi_k)
    :=
    \mathbb{E}_{s\sim d^{\pi_k},\,a\sim\pi_k(\cdot|s)}
    \left[
        \frac{\pi_{k+1}(a\mid s)}{\pi_k(a\mid s)}
        A^{\pi_k}(s,a)
    \right].
    \label{eq:ratio_surrogate}
\end{equation}
This local surrogate, denoted by $\tilde{\Delta}$, forms the basis of practical proximal policy optimization algorithms. 
TRPO maximizes such a surrogate subject to a KL-based trust-region constraint between the old and updated policies, while Proximal Policy Optimization (PPO)~\citep{SchulmanWDRK17PPO} replaces the constrained optimization with a clipped probability-ratio objective that discourages excessively large policy updates. Based on PPO, Group Relative Policy Optimization (GRPO)~\citep{Shao24GRPO} adopts a group-based advantage estimation which is free of learning the value function.
\vspace{-0.2cm}
\paragraph{Training-Inference Mismatch in LLM RL}
As discussed above, modern LLM RL pipelines may induce different action distributions in the training and inference engines, even when they share synchronized model parameters.
We model this training-inference mismatch by distinguishing the training policy $\textcolor{blue}{\pi}$ from the inference policy $\textcolor{red}{\mu}$: at update step $k$, with the same parameters, $\textcolor{blue}{\pi_k}$ is the policy represented inside the training engine, whereas $\textcolor{red}{\mu_k}$ is the policy actually used for rollout generation and deployment in the inference engine. 
The training-inference mismatch manifests as $\textcolor{blue}{\pi}\neq\textcolor{red}{\mu}$, even though they share the same model parameters.
This mismatch thus results in new learning issues regarding non-stationarity and suboptimality different from the notorious ones in traditional RL~\citep{Hasselt18DRLDeadly,TangB24CHAIN,Wu2026SWD}.
In the following of this paper, we will use the notations to explicitly highlight the distinction between training policies and inference policies.
This distinction provides the basis for our method, where we revisit policy improvement from the perspective of the inference policy rather than the training-side surrogate.

